\subsection{Strict Laurent positivity and the source's changing cells}
The source's rank-four two-row basis has a feature that an
effective graded matrix kernel cannot preserve: one generic cell
splits into two specialized cells at \(v=1\).
Here the required basis subsets, rather than the abstract
existence of a specialized decomposition, supply the obstruction.

\begin{finding}[\fid{F58}: source-cell compatibility forces cancellation]
\label{f:cellpositivityspecialization}
Let \(T_1,\ldots,T_r\) be matrices with entries in
\(\Q_{\ge0}[v,v^{-1}]\), acting in a fixed finite basis.
At every positive real specialization \(v=a>0\), the generator
support graph, the all-algebra cell preorder, and the coordinate
invariant subspaces are unchanged. The same assertion holds for
every cellular coordinate subquotient.

Consequently the actual six-dimensional two-row module
\[
N_6=\check M_{\{(3,1),(2,2)\}}
\]
has no regular specializing basis which both makes every named
\(P_i\) matrix a nonnegative Laurent polynomial matrix and makes
its source-prescribed three-dimensional specialized
\(M_{(2,1,1)}\) submodule a basis-subset span.
The same conclusion holds with the \(Q_i\) in place of the \(P_i\).

More generally, no rank-four tensor basis satisfying the generic
cell and specialized quasi-cell requirements of IV,
Conjecture~19.1(ii)--(iii), can additionally have all its named
\(P_i\) matrices in \(\Q_{\ge0}[v,v^{-1}]\).
The \(Q_i\)-only version is equally impossible.
These are additional strict positivity requirements, not claims
made by Conjecture~19.1 as stated.
\end{finding}

\begin{proof}
\par\smallskip\noindent\textbf{Positive specialization reflects zero.}\quad
More generally let \(k\subseteq\mathbb R\) be an ordered field,
\(A=k[v,v^{-1}]\), and \(K=k(v)\).
For a finite Laurent polynomial
\(f=\sum_e c_ev^e\) with \(c_e\ge0\), and for \(a>0\),
\[
f(a)=\sum_e c_ea^e=0
\quad\Longleftrightarrow\quad
c_e=0\text{ for every }e
\quad\Longleftrightarrow\quad
f=0.
\tag{C1}
\]
Every nonzero summand is strictly positive.
Thus the zero pattern of each named generator matrix is
unchanged at a positive specialization.

Write matrices on the right of row vectors.
For a subset \(S\) of the basis indices, its coordinate span
is invariant under all generators exactly when
\[
(T_i)_{st}=0
\qquad
\text{for every }i,\ s\in S,\ t\notin S.
\tag{C2}
\]
Condition (C1) makes these equations equivalent before and
after specialization. A coefficient which is zero over \(K\)
is identically zero in \(A\).

\par\smallskip\noindent\textbf{The source cell preorder.}\quad
IV, \S2.4, pp.~18--19, equation~(17), defines the cell preorder
using nonzero coefficients of arbitrary algebra elements.
Its cellular submodules are explicitly spans of subsets
of the designated basis. A cellular subquotient is
\(A\Gamma'/A\Gamma''\), where
\(\Gamma''\subseteq\Gamma'\) span coordinate submodules,
with quotient basis \(\Gamma'\setminus\Gamma''\).
The right-action version is stated on p.~19 \cite{IV}.

The all-algebra preorder agrees with reachability in the
directed graph of named generator entries.
Every generator edge is an allowed algebra-action relation.
Conversely, an element of the generated algebra is a scalar
linear combination of words. If one of its matrix entries is
nonzero, some word has a nonzero entry in that position.
Expanding that word entry as a sum over intermediate states
produces at least one nonzero product, hence a path of nonzero
generator entries. Taking transitive closures proves the
claim, with the source's orientation for the preorder.
This argument itself does not require positivity.

By (C1) the generator graph is the same generically and at
\(a>0\). Its reachability relation, strongly connected
components, and corresponding coordinate invariant subsets
are therefore the same. In particular a generic cell cannot
split into specialized cells in this basis.

For a cellular subquotient with basis
\(\Gamma'\setminus\Gamma''\), the induced matrices are exactly
the corresponding submatrices of the original matrices.
They retain nonnegative Laurent coefficients.
The coordinate modules and quotient are free, so their
specialization commutes with the quotient.
The preceding conclusions apply to this subquotient as well.

\par\smallskip\noindent\textbf{The actual six-dimensional cell.}\quad
IV, Theorem~13.4, pp.~74--75, includes
\(N_6=\check M_{\{(3,1),(2,2)\}}\) as a generic simple of the
two-row nonstandard quotient.
It occurs in the actual tensor representation already when
\(\dim V=\dim W=2\).
Example~19.3, pp.~133--135, constructs its designated cell,
drawn in Figure~21 on p.~135.
The text on p.~134 explicitly states that the top row spans
a specialized \(\Q S_4\)-cellular submodule isomorphic to
\(M_{(2,1,1)}\), while the bottom row gives the cellular
quotient \(M_{(3,1)}\).
Both dimensions are three \cite{IV}.

The graph displays the exact mechanism.
Order the six states as
\[
(v_{12},v_{13},v_{23},v_3,v_2,v_1),
\]
with the three top-row states first.
Put
\[
\delta=v+v^{-1},\qquad D=\delta^2,\qquad
a_2=D-2,\qquad c_0=D-4=(v-v^{-1})^2.
\]
Its only top-to-bottom arrow is
\[
v_{13}\xrightarrow{\,Q_2,\ c_0\,}v_2.
\tag{C3}
\]
The three top-row states are mutually reachable, as are
the three bottom-row states.
Bottom-to-top arrows remain nonzero at \(v=1\).
Generically \(c_0\ne0\), so the entire six-state graph is
strongly connected.
At \(v=1\), one has \(D=4,\ a_2=2,\ c_0=0\).
The arrow (C3) disappears and the top-row coordinate span
becomes an invariant subspace.

The complementary generators satisfy \(P_i=DI-Q_i\).
Their off-diagonal supports are therefore identical,
and they generate the same unital algebra.
On the actual specialized module,
\[
P_i(1)=2(1+s_i),\qquad Q_i(1)=2(1-s_i),
\tag{C4}
\]
by the rank-one specialization in IV, \S11.1, p.~59.
Thus either set of specialized generators generates the
ordinary diagonal symmetric-group image.
The exponent convention issue recorded in \fid{F57} does
not affect (C4), since both preliminary conventions agree
at \(v=1\).

For full precision, a regular basis here can be a basis over
\[
R=\Q[v,v^{-1}]_{(v-1)}
\]
of a free form of the actual generic module, with a basis
after specialization at \(1\).
If its action entries lie in \(\Q_{\ge0}[v,v^{-1}]\),
its Laurent span is a stable free \(A\)-module having the
same generic action and the same labeled specialized basis.
No equality with an unspecified global integral form is
being assumed.

Now suppose the three-dimensional specialized submodule
is spanned by a subset of such a positive basis.
Condition (C2) lifts that subset to a generic invariant
three-dimensional subspace, contradicting the generic
simplicity of \(N_6\).
This proves the assertion for either generator set.

\par\smallskip\noindent\textbf{The tensor-basis strengthening.}\quad
Conjecture~19.1(ii) specifies the generic nonstandard
cellular subquotients and their right modules.
Condition (iii) specifies the specialized cells, or
quasi-cells, as a refinement inside each generic cell,
with their prescribed Specht modules.
For \(\alpha=\{(3,1),(2,2)\}\), this refinement is the
two three-dimensional quasi-cells just described.

If all named \(P_i\) matrices of such a tensor basis
were coefficientwise nonnegative Laurent polynomials,
the induced matrices on this six-dimensional cellular
subquotient would have the same property.
Its cell partition would then be unchanged at \(v=1\),
contrary to condition (iii).
Equivalently, the global cell preorder cannot change
under positive specialization, whereas these two source
conditions require a nontrivial refinement.
The identical argument applies to the \(Q_i\).
\end{proof}

\paragraph{Why an arbitrary positive basis is not excluded.}
The coordinate premise is essential.
For example the two matrices
\[
A(v)=\begin{pmatrix}1&1\\1&1\end{pmatrix},
\qquad
B(v)=\begin{pmatrix}1&v\\1&1\end{pmatrix}
\tag{C5}
\]
have entries in \(\mathbb N[v]\).
The only invariant row lines for \(A\), even over an
algebraic closure of \(\Q(v)\), are those generated by
\((1,1)\) and \((1,-1)\), because its eigenvalues are
\(2\) and \(0\).
But
\[
(1,1)B=(2,v+1),\qquad
(1,-1)B=(0,v-1),
\]
so neither line is \(B\)-invariant for \(v\ne1\).
The generic generated algebra therefore acts absolutely
irreducibly.
At \(v=1\), the matrices coincide and the module splits
into those two eigenspaces.
Neither eigenspace is a subset span of the standard basis.
The support graph remains complete and has one cell.
This elementary example is not asserted to satisfy the
nonstandard Hecke relations.

Thus \fid{F58} does not exclude an arbitrary positive
basis of the full tensor module, or of \(N_6\), whose
specialized classical summands are non-coordinate.
It also does not exclude coefficients nonnegative merely
as functions on \(v>0\):
\(c_0=(v-v^{-1})^2\) is such a function, but has Laurent
coefficient \(-2\) and vanishes at \(1\).
Signed/factored coefficients and virtual or derived
coefficients can supply precisely this cancellation.
The finding is a restriction on an added strict
effective-kernel comparison with the source's designated
cell structure, not a contradiction to the source's
already constructed rank-four basis.

% Separate paragraph for Part III.
\paragraph{\fid{F58}'s specialization in Grothendieck's language.}
Finite graded coefficient spaces are representations
of the split torus \(\mathbb G_m\).
Their effective characters lie in
\(\mathbb N[v,v^{-1}]\), and forgetting the grading gives
the dimension augmentation
\[
\epsilon:\mathbb Z[v,v^{-1}]\longrightarrow\mathbb Z,
\qquad v\longmapsto1.
\]
This map reflects zero on the effective cone:
a nonzero finite direct sum of weight spaces has positive
dimension.
The representation framework is SGA~3,
I, 4.7.2--3 and cote~90, scan~17
\cite{SGA3,Archive}.
The representation and \(\lambda/\sigma\) vocabulary is
SGA~6, V, \S\S2--3 and cote~133, scan~18
\cite{SGA6,Archive};
no \(\lambda\)-ring structure on the noncommutative
word algebra is claimed.

In affine language the chosen basis gives a trivial
bundle on \(\Spec A\), and a basis subset gives a
direct-summand subbundle.
Its invariance is the vanishing of the cross-block
matrix entries.
The module/sheaf dictionary is EGA~I,
Theorem~1.3.7, p.~86, Corollary~1.3.8, p.~87,
and Theorem~1.4.1, pp.~90--92 \cite{EGA}.
Once the graded-character parameter is identified
with the source parameter, effectiveness prevents
these cross-block equations from vanishing only at
the positive identity fiber.
The actual class
\[
c_0=v^2-2+v^{-2}
\]
is in \(\ker\epsilon\) but is not effective.
Ordinary representation \(K_0\) therefore permits
the required edge cancellation; an effective graded
matrix kernel in the prescribed cell basis does not.
A virtual or derived replacement still needs its
own designated basis and arithmetic comparison.

